\textbf{Memorisation.} Zhang et al.~\cite{zhang2016understanding} showed that standard networks fit random labels, so capacity alone does not prevent memorisation. Arpit et al.~\cite{arpit2017closer} found that networks fit simple, shared patterns first and noisy examples later, and that regularisation and data affect this differently. Patel and Sastry~\cite{patel2021memorization} study how the choice of loss function changes memorisation of noisy labels, which is close to our question but with different losses.

\textbf{Sample selection.} MentorNet~\cite{jiang2017mentornet} learns a curriculum that down-weights probable outliers; Co-teaching~\cite{han2018co} trains two networks that each pass their small-loss examples to the other, with a forget rate that follows the noise rate and increases over training; DivideMix~\cite{li2020dividemix} splits data into clean and noisy parts with a mixture model and continues semi-supervised. Our selection baseline is the one-network, per-batch variant of the small-loss rule; it is a reimplementation and omits the cross-update that gives Co-teaching its name, so it should not be read as a reproduction of that method.

\textbf{Regularisation.} Early-learning regularisation (ELR)~\cite{liu2020early} prevents memorisation by steering towards early-training predictions. Mixup~\cite{zhang2017mixup} trains on convex combinations of examples and labels and was reported to reduce memorisation of corrupted labels. Label smoothing~\cite{mller2019when} replaces hard targets by a mixture with the uniform distribution. Lukasik et al.~\cite{lukasik2020does} analyse whether it mitigates label noise, relating it to a form of loss correction, while Wei et al.~\cite{wei2021smooth} report that its advantage vanishes in a high-noise regime. Our experiments test these two statements on a tiny dataset.
